% ECONOMICS_PROBLEM: {"id": "C63-35", "title": "Scalable heterogeneous-agent macroeconomics", "jel_code": "C63", "source_id": "OP-0016"}
% !TeX program = xelatex
\documentclass[11pt,a4paper]{article}
\usepackage[margin=23mm]{geometry}
\usepackage{fontspec}
\setmainfont{Latin Modern Roman}
\usepackage{amsmath,amssymb,mathtools}
\usepackage{xcolor,enumitem,booktabs,longtable,array,xurl,hyperref}
\definecolor{muted}{HTML}{53636C}
\hypersetup{colorlinks=true,urlcolor=blue,linkcolor=blue}
\setlength{\parindent}{0pt}
\setlength{\parskip}{5pt}
\newcommand{\R}{\mathbb R}
\newcommand{\E}{\mathbb E}
\newcommand{\N}{\mathbb N}
\newcommand{\ind}{\mathbf 1}
\newcommand{\Var}{\operatorname{Var}}
\newcommand{\Cov}{\operatorname{Cov}}
\newcommand{\supp}{\operatorname{supp}}
\DeclareMathOperator*{\argmin}{arg\,min}
\DeclareMathOperator*{\argmax}{arg\,max}
\newcommand{\entrymeta}[3]{{\sffamily\small\color{muted}\textbf{\texttt{#1}}\quad|\quad #2\par #3\par}}
\newcommand{\Problemref}[1]{\texttt{#1}}
\newcommand{\JELref}[2]{\texttt{#2} (JEL \texttt{#1})}
\begin{document}
\section*{Scalable heterogeneous-agent macroeconomics}
\textbf{Problem ID:} \texttt{C63-35}\quad \textbf{JEL:} \texttt{C63}\par
% BEGIN ECONOMICS STATEMENT
\label{problem:OP-0016}
\label{atlas:prob:A6}

\noindent\textbf{Original atlas ID:} A6. \textbf{Formulation:} editorial specification of a research family.

\subsection*{Definitions and assumptions}

Specify a heterogeneous-agent economy with idiosyncratic state $s\in S$, assets $a$ in a bounded interval with a borrowing limit, aggregate shock $z\in Z$, and a population probability measure $\mu\in\mathcal P(S\times A)$. The aggregate state is $(z,\mu)$; prices, household policies, and market-clearing conditions are equilibrium objects. Individual policies must solve the stated discounted optimization problems, and $\mu'$ must equal the distribution induced by those policies and the transition kernels. Let $\mathcal D$ be a declared region of aggregate states and $f$ collect equilibrium policy, price, and distribution-transition functions. An operator $\mathcal R(f)$ records equality residuals for feasibility, market clearing, and distributional evolution, positive parts of violated optimality inequalities, and complementarity errors at binding constraints; satisfied slack inequalities must contribute zero violation. Its norm, discretization, and treatment of binding constraints are specified.

\subsection*{Formal research question}

For a restricted primitive class $\mathcal H$, can one compute $\hat f$ with a certified bound
\[
\sup_{h\in\mathcal H}\|\mathcal R_h(\hat f_h)\|_{\mathcal D}\le\varepsilon
\]
at a feasible computational cost, including states away from a steady state and transitions across binding constraints? Identify stability conditions under which the residual controls distance to an equilibrium:
\[
\operatorname{dist}(\hat f_h,\mathcal E_h)\le C_h\|\mathcal R_h(\hat f_h)\|_{\mathcal D}.
\]
For estimation, propagate numerical error into moments, likelihoods, and counterfactual uncertainty, so solver bias is distinguishable from sampling error. Benchmark methods on identical primitives, initial distributions, shocks, hardware budgets, and accuracy criteria; local linear accuracy and global state-domain accuracy must be reported separately.

\subsection*{Interpretation and scope}

The residual-to-error bound is a target requiring an explicit stability theorem; a small residual alone does not imply proximity to the correct equilibrium when the operator is ill conditioned or equilibria are multiple. Household Bellman contraction with fixed prices does not establish contraction of the coupled market-clearing economy. Asset truncation, distribution compression, quadrature, and tail losses all contribute numerical error. ``Global'' must name a state domain and norm, since an arbitrary probability-measure state space cannot be exhaustively validated by a few simulated paths. Occasionally binding inequalities require complementary optimality checks. A useful benchmark also reports how approximation error changes estimated parameters and distributional policy conclusions.

\subsection*{Source audit and status}

Aiyagari (1994), Krusell--Smith (1998), and Kaplan--Moll--Violante (2018) are verified. Original link [18] names the 2016 NBER working-paper version of HANK, whereas the cited 2018 journal version is used here. Auclert--Bardóczy--Rognlie--Straub (2021), added as [S4], supplies an important efficient sequence-space solution and estimation method, preventing the impression that the computational field lacks established tools. The remaining question concerns verified nonlinear/global accuracy for explicitly rich models and estimation tasks. The operator and uniform guarantee above are editorial numerical-analysis specifications, not a claim that all heterogeneous-agent equilibria are unresolved or that local methods necessarily fail.

\subsection*{References}

\begin{enumerate}

\item[S1] S. Rao Aiyagari (1994). \emph{Uninsured Idiosyncratic Risk and Aggregate Saving}. Quarterly Journal of Economics 109(3), 659--684. \url{https://academic.oup.com/qje/article-abstract/109/3/659/1838287}. \textit{Locator:} abstract; incomplete-markets equilibrium. \textit{Establishes:} A heterogeneous-agent saving and wealth-distribution benchmark. \textit{Does not establish:} a globally accurate solver for arbitrary aggregate shocks.

\item[S2] Per Krusell and Anthony A. Smith, Jr. (1998). \emph{Income and Wealth Heterogeneity in the Macroeconomy}. Journal of Political Economy 106(5), 867--896. \url{https://www.journals.uchicago.edu/doi/10.1086/250034}. \textit{Locator:} abstract; approximate aggregation computation. \textit{Establishes:} A practical computation with aggregate risk and heterogeneous households. \textit{Does not establish:} a model-independent global error guarantee.

\item[S3] Greg Kaplan, Benjamin Moll, and Giovanni L. Violante (2018). \emph{Monetary Policy According to HANK}. American Economic Review 108(3), 697--743. \url{https://pubs.aeaweb.org/doi/10.1257/aer.20160042}. \textit{Locator:} abstract; heterogeneous-agent monetary model. \textit{Establishes:} A distribution-rich account of monetary transmission. \textit{Does not establish:} global nonlinear estimation for every heterogeneous-agent economy.

\item[S4] Adrien Auclert, Bence Bardóczy, Matthew Rognlie, and Ludwig Straub (2021). \emph{Using the Sequence-Space Jacobian to Solve and Estimate Heterogeneous-Agent Models}. Econometrica 89(5), 2375--2408. \url{https://onlinelibrary.wiley.com/doi/10.3982/ECTA17434}. \textit{Locator:} abstract; sequence-space Jacobian method. \textit{Establishes:} Efficient solution and estimation using sequence-space derivatives. \textit{Does not establish:} uniform global accuracy under all large shocks and binding constraints.

\end{enumerate}

\subsection*{Common conventions: Source mathematical conventions}

Unless an entry states otherwise, agent and firm sets are finite. State and action spaces are measurable, strategies and outcome maps are measurable, probability laws are specified, and expectations exist for the targets under discussion. Infinite-horizon discount factors lie in $(0,1)$ and discounted objectives are finite. Norms, tolerances and complexity input sizes must be specified for computational comparisons. Constants or primitives that appear locally are inputs, not empirical facts asserted by this catalogue.

\textbf{Identification.} A statistical model $m\in\mathcal M$ implies an observable law $P_m$ and target $\theta(m)$. At law $P$, its identified set is
\[
\Theta(P;\mathcal M)=\{\theta(m):m\in\mathcal M,\ P_m=P\}.
\]
Point identification holds when this set is a singleton, or equivalently when $P_{m_1}=P_{m_2}$ implies $\theta(m_1)=\theta(m_2)$ for all compatible models. Sharp bounds mean bounds attained, or approached when the boundary is not attained, by compatible models. Writing this set is a definition, not a solution: the substantive task is to establish informative restrictions and derive or estimate the set. An unrestricted-confounding model may give only trivial bounds.

\textbf{Inference.} Identification, estimability and valid uncertainty quantification are separate questions. Fix $\alpha\in(0,1)$. A confidence procedure $C_n$ over a declared law class $\mathcal P$ has asymptotic uniform coverage at level $1-\alpha$ when
\[
\liminf_{n\to\infty}\inf_{P\in\mathcal P}\Pr_P\{\theta(P)\in C_n\}\geq1-\alpha.
\]
For set-identified targets the coverage event must be defined separately, for example coverage of the full identified set. If an entry uses identification-indexed models instead of uniquely indexed laws, the same coverage convention is applied over those models. Pointwise limits do not establish uniform validity, and asymptotic validity does not guarantee finite-sample precision. Sample sizes, dependence structures and weak-identification sequences are part of each application.

\textbf{Counterfactuals.} $Y_i(a)$ is a potential outcome under a specified intervention, including its timing, implementation and versions. It is not $Y_i$ conditional on voluntarily choosing $a$. A full intervention vector permits interference; shorthand scalar treatment notation does not silently impose no interference. Assignment, selection, attrition, anticipation and measurement must be specified. A claim about an instrument's local effect is not automatically a population policy effect.

\textbf{Economic equilibrium.} A counterfactual model includes best responses, constraints, feasibility, market clearing, state transitions and expectations consistent with the stated information. When several equilibria exist, either a declared selector is an input or the target is set-valued. Comparisons across policy regimes must use the stated selection convention. Reduced-form relationships are not presumed invariant after an equilibrium-changing intervention.

\textbf{Optimization and welfare.} Optimization is over the stated feasible set. If attainment is not established, $\sup$ or $\inf$ and $\varepsilon$-optimal choices replace an asserted maximizer or minimizer. For example, with value $V=\sup_{a\in\mathcal A}W(a)<\infty$, an $\varepsilon$-optimal set is $\{a\in\mathcal A:W(a)\geq V-\varepsilon\}$ for $\varepsilon>0$. Benefits, costs, risk and health or environmental outcomes can be aggregated only using declared units and conversion weights. Interpersonal utility weights, the treatment of fiscal transfers and foreign profits, discounting, and the behavioral welfare criterion are explicit inputs. A comparative-static derivative additionally requires existence and appropriate differentiability of the selected counterfactual.

\textbf{Sources.} Local labels [S1], [S2], and so forth restart in every question. Locators identify the relevant abstract, model, section or result at the level actually checked; a bibliographic or abstract-level verification is not represented as inspection of an entire proof. Working papers are distinguished from later publications and changing versions. Source summaries are paraphrased. All URL checks and editorial formulations refer to the audit date stated above.
% END ECONOMICS STATEMENT
\end{document}
